% Input fragment: refinement of the shared-point construction.
\subsection{Comparison order fifty-two with an explicit larger anchor}
\label{subsec:shared-precision52}

The high comparison order can be lowered from $64$ to $52$ while
keeping beta order $16$ and scalar-height order $48$.  This does not
change the number of retained coordinates of the shared construction.
It shortens eight reciprocal logarithms, and can therefore reduce the
number of variables in a quartic circuit presentation.  Such a circuit
count is a separate, implementation-dependent assertion.

\begin{proposition}\label{prop:shared-precision52}
Assume the arithmetic inputs of the shared-point construction and
$1\le H<2^{4783227}$.  Replace every order-$64$ comparison by order
$52$, and replace the corresponding logarithms $zG_{15}(Df^2)$ by
$zG_{12}(Df^2)$.  Retain the beta order $16$, the scalar order $48$,
and all other equations and guards.  Either of the following size
conditions is sufficient:
\[
 |a|>(21H)^{13/2},
\]
or the semilocal anchor with fixed constant
\[
 C=3^{2^{26}},\qquad
 u^2+6v^2-2w^2-3(f_0^2+f_0f_1-f_1^2)^2+2C^4=0.
\]
Here $f_0+f_1\sqrt2=f(S)$.  Alternatively the finite exclusion
condition can use $M_{52}=2^{31091008}$.  Each version has the same
integer input projection as the order-$64$ construction under the
stated arithmetic inputs.
\end{proposition}

\begin{proof}
We first check that the changed logarithm has the required properties.
The formal identity
\[
 \log_E(z)=z\sum_{j\ge0}c_j(Df^2)^j,
 \qquad c_1=-2/5,\quad
 c_j=c_{j-1}(4j-5)/(4j+1),
\]
is independent of the chosen truncation.  At a place above
$p>76176$, normalized by $v(p)=1$, a nonzero point in the formal
group has $r=v(z)\ge1$ and $v(Df^2)=4r$.  The denominator of $c_j$
divides $(4j+1)!$.  For $j\ge13$ with $4j+1<p$, the $j$th term
therefore has valuation at least $52r$.  Otherwise $j\ge19044$ and
\[
 v\bigl(c_j(Df^2)^j\bigr)
 \ge4jr-\frac{4j+1}{p-1}
 \ge4jr-j\ge52r.
\]
These lower bounds tend to infinity, so they control the whole tail.
Consequently
\[
 v\bigl(\log_E(z)-zG_{12}(Df^2)\bigr)\ge53r,
 \qquad
 \frac{zG_{12}(Df^2)}{\log_E(z)}\in1+p^{52}\mathcal O_{K,\mathfrak p}.
\]
This proof applies to either a split or an inert good prime, since
both local extensions are unramified.  The source's global
nonvanishing argument applies to every $G_J$ with $J\ge2$ and hence
to $G_{12}$.  Thus all logarithm denominators remain nonzero.

The two-anchor comparison and the detector proof use only the
specified precision; they now give order $52$ for the three Lucas
labels and order $16$ for beta.  The transformed contact-height
estimate is unchanged.  Put $a_0=|a|$ and
$B=\max(|n|,|a|,|t|,|b|)$.  If the integer Lucas residual
$\Delta$ is nonzero, the resulting estimates are
\[
 B^2-6<7H|\Delta|^{24/52},\qquad |\Delta|<5B^4.
\]
Since $B,H\ge1$, they imply
\[
 B^{2/13}<\bigl(7\,5^{6/13}+6\bigr)H<21H,
 \qquad B<(21H)^{13/2}.
\]
The strict constant inequality follows by raising
$5^{6/13}<15/7$ to the thirteenth power; equivalently,
$5^6 7^{13}<15^{13}$.  In particular replacing the factor by
$14H$ would not be justified: $7\,5^{6/13}>14$.
The size hypothesis contradicts $B\ge a_0$, so $\Delta=0$ and
the unchanged Lucas lemma gives $\log B>a_0-3$.

The scalar estimate remains $h(\beta)\le H\sqrt{a_0}$.
If $\beta\ne n$, the order-$16$ argument remains
\[
 \log B<2\log(32H)+3h(\beta).
\]
The new size hypothesis implies $a_0>(20H)^2$, because
$21^{13}>20^4$ and $H\ge1$.  Hence
$h(\beta)<a_0/20$.  Also
$\log(32H)<20H^2<a_0/20$ for $H\ge1$.
It follows that $\log B<a_0/4<a_0-3$, a contradiction.  Thus
$\beta=n$ and $X\in\mathbb Z$ as before.

For the explicit constants,
\[
 \log_2\bigl((21H)^{13/2}\bigr)
 <\frac{13}{2}(5+4783227)=31091008.
\]
The exact source subgroup calculation for an anchor
$C=3^{2^q}$ gives
\[
 0\ne a\in2\,3^{2^{q-1}-1}\mathbb Z.
\]
For $q=26$, using $3^2>2^3$,
\[
 \log_2|a|\ge1+(2^{25}-1)\log_2 3
 >1+\frac32(2^{25}-1)=50331647.5>31091008.
\]
The semilocal anchor is therefore sufficient.  The alternative
indexed exclusion with $M_{52}=2^{31091008}$ is sufficient by the
same inequality.  Its enormous number of factors is part of the
explicit indexed specification, not an expanded coefficient table.

Completeness still chooses the anchor index in an arbitrarily
divisible nonzero subgroup, or larger than a fixed finite exclusion.
All the prime-pattern, congruence, local-square, and real-smallness
arguments are unchanged: lowering a required comparison precision
only weakens the necessary formal divisibilities.  This proves the
claimed equality of input projections.
\end{proof}

For the detector arithmetic, with $F=F_5(A\lambda)$ and
$\delta=A\lambda-1$, the three possible weights become
\[
 q_1=\frac{A^5\delta^2}{F},\qquad
 q_{16}=\frac{A^{155}\delta^{152}}{F^{31}},\qquad
 q_{52}=\frac{A^{515}\delta^{512}}{F^{103}}.
\]
These follow from $q_m=A^{10m-5}\delta^{10m-8}/F^{2m-1}$.
There are eight order-$52$ logarithms and two order-$16$
logarithms.  Each changed logarithm loses exactly three terms.
The fixed coefficient in the anchor is $2C^4=2\cdot3^{2^{28}}$;
using it as a literal coefficient adds no variable.
